%! Piecewise Linear Learning Functions — Unit 8, Lesson 8.1 — Exit Ticket (Answer Key)
\documentclass[10pt]{article}
\usepackage{linalg-article}
\usepackage{linalg-key}   % pulls in -boxes; do NOT also load -boxes
\newcommand{\TallMath}[1]{$\displaystyle #1\rule[-1.4em]{0pt}{3.2em}$}
\newcommand{\relu}{\mathrm{ReLU}}

\begin{document}

\pageheader{Unit 8, Lesson 8.1}{Exit Ticket --- Answer Key}
\namedateperiod

\vspace{0.15in}

No notes. Show your reasoning.

\begin{enumerate}[leftmargin=1.8em, itemsep=16pt]
  \item \textbf{Build the shape.} Let $F(x)=\relu(x)-2\,\relu(x-2)+\relu(x-4)$. Compute
        $F(0)$, $F(2)$, $F(4)$, and $F(6)$, and describe the shape the graph makes.
        \ansline{$F(0)=0$; $F(2)=2$; $F(4)=4-2(2)=0$; $F(6)=6-2(4)+2=0$. The graph rises with slope $1$ to the peak $(2,2)$, falls with slope $-1$ back to $0$ at $x=4$, then stays flat --- a \textbf{tent} (triangle).}
  \item \textbf{Folds and pieces.} A piecewise-linear function is built from $4$ ReLUs. How many folds does
        it have, and how many straight pieces?
        \ansline{$4$ folds and $4+1=5$ straight pieces.}
  \item \textbf{Explain.} In one sentence, why can a sum of ReLUs bend to fit rising-and-falling data when a
        single straight line cannot?
        \ansline{Each ReLU adds a fold where the slope can change, so a sum of them can turn up and then down; a single line has no folds and one fixed slope, so it can only go straight.}
\end{enumerate}

\begin{teachernote}
Sort responses into ``got it / arithmetic slip / concept slip.'' Item~1 checks building a piecewise-linear
function: turn on only the ReLUs whose fold is left of $x$ (peak $(2,2)$, back to $0$ at $x=4$, then flat).
Watch for students who forget a ReLU only switches on past its fold, or who miscompute $-2\,\relu(x-2)$.
Item~2 is the counting rule $N\to N+1$. Item~3 is the target idea: \emph{folds are what let a graph change
direction}. If many stumble on item~1's shape or item~3, reopen next lesson with a 30-second ``each ReLU $=$
one fold; folds let the graph turn'' recap.
\end{teachernote}

\end{document}
